\section*{第 \S\arabic{exercisegroup} 节习题}
\phantomsection
\addcontentsline{toc}{section}{第 \protect\S\arabic{exercisegroup} 节习题}

\begin{enumerate}
\def\labelenumi{\arabic{enumi})}
\tightlist
\item
  \begin{enumerate}
  \def\labelenumii{\alph{enumii})}
  \tightlist
  \item
    沿用 IV, p.~165 的记号，假设 \(\mathbf{f}\) 在 \(\mathrm{J} \times \mathrm{S}\) 上一致连续。不用选择公理证明 IV, p.~166 的命题 3。（设 \(\delta\) 使得关系 \(|t_2 - t_1| \leqslant \delta\) 、\(\| \mathbf{x}_2 - \mathbf{x}_1 \| \leqslant \delta\) 蕴含 \(\| \mathbf{f}(t_2, \mathbf{x}_2) - \mathbf{f}(t_1, \mathbf{x}_1) \| \leqslant \varepsilon\) ；把 J 分成长度 \(\leqslant \inf(\delta, \delta / M)\) 的区间，并用归纳法定义近似解。）
  \end{enumerate}
\end{enumerate}

\begin{enumerate}
\def\labelenumi{\alph{enumi})}
\setcounter{enumi}{1}
\tightlist
\item
  当 E 是有限维且 f 在 \(I \times H\) 上是利普希茨的时，不用选择公理证明 IV, p.~166 的命题 3。（注意，对每个 \(\delta > 0\) ，存在 S 的有限多个点 \(x_{i}\) ，使得 S 的每个点都以 \(\leqslant \delta\) 的距离靠近某个 \(x_{i}\) ；然后照 a) 的做法进行，考虑规则函数 \(t \mapsto \mathbf{f}(t, \mathbf{x}_{i})\) ，它们的个数是有限的。）
\end{enumerate}

\begin{enumerate}
\def\labelenumi{\arabic{enumi})}
\setcounter{enumi}{1}
\item
  给定两个数 \(b > 0\) 、\(M > 0\) 以及任意一个数 \(\varepsilon > 0\) ，给出一个数量微分方程 \(x' = g(x)\) 的例子，使得 \(|g(x)| \leqslant M\) （对 \(|x| \leqslant b\) ），并且它有一个积分 \(x = u(t)\) ，该积分在区间 \(] - b/M - \varepsilon, b/M + \varepsilon[\) 上连续，但在点 \(x = \frac{b}{M} + \varepsilon\) 处没有有限极限（这样定义 \(g\) ：使所考虑的积分在 \(] - b/M - \varepsilon, b/M + \varepsilon[\) 上有连续导数，且该导数等于常数 \(M\) 于 \(] - b/M, b/M[\) 上）。
\item
  设 \(S \subset H\) 是以 \(\mathbf{x}_0\) 为中心、\(r\) 为半径的开球，\(\mathbf{f}\) 是以常数 \(k > 0\) 为利普希茨常数的、在 \(I \times S\) 上的利普希茨函数；再设 \(t \mapsto \mathbf{f}(t, \mathbf{x}_0)\) 在 \(I\) 上有界，并以 \(M_0\) 表示 \(\| \mathbf{f}(t, \mathbf{x}_0) \|\) （对 \(t \in I\) ）的上确界。证明：对所有的 \(t_0 \in I\) ，存在一个积分 \(\mathbf{u}\) ——\(\mathbf{x}' = \mathbf{f}(t, \mathbf{x})\) 的积分——它取值于 \(S\) ，等于 \(\mathbf{x}_0\) （在点 \(t_0\) 处），并且定义在 \(I\) 与区间 \(]t_0 - \lambda, t_0 + \lambda[\) 的交上，其中
\end{enumerate}

\[
\lambda = \frac {1}{k} \log \left(1 + \frac {k z}{M _ {0}}\right)
\]

（注意，有 \(\| \mathbf{u}(t) - \mathbf{x}_0\| \leqslant \mathbf{M}(t - t_0) + k\int_{t_0}^t\| \mathbf{u}(s) - \mathbf{x}_0\| ds\) （对 \(t > t_0\) ））。

* 4) a) 设 \(I = ]t_0 - a, t_0 + a[\) 是 \(\mathbf{R}\) 中的开区间，\(S\) 是以 \(\mathbf{x}_0\) 为中心、\(r\) 为半径的开球（在 \(E\) 中），\(\mathbf{f}\) 是 \(I \times S\) 上的局部利普希茨函数。设 \((s, z) \mapsto h(s, z)\) 是一个 \(\geqslant 0\) 函数，对 \(0 \leqslant s \leqslant a\) 与 \(0 \leqslant z \leqslant r\) 有定义且连续，使得对每个 \(s \in [0, a]\) ，映射 \(z \mapsto h(s, z)\) 是递增的；假设 \(\| \mathbf{f}(t, \mathbf{x}) \| \leqslant h(|t - t_0|, \| \mathbf{x} - \mathbf{x}_0 \|)\) 在 \(I \times S\) 上成立。设 \(\varphi\) 是区间 \([0, \alpha]\) （\(\alpha < a\) ）上一个规则函数的原函数，取值于 \([0, r[\) ，使得 \(\varphi(0) = 0\) 且 \(\varphi'(s) > h(s, \varphi(s))\) （在 \([0, \alpha]\) 上，除一个可数集外）。证明积分 \(\mathbf{u}\) ——\(\mathbf{x}' = \mathbf{f}(t, \mathbf{x})\) 的、等于 \(\mathbf{x}_0\) 于点 \(t_0\) 的积分——定义在 \(]t_0 - \alpha, t_0 + \alpha[\) 上，并且在这个区间上 \(\| \mathbf{u}(t) - \mathbf{x}_0 \| \leqslant \varphi(|t - t_0|)\) 。

\begin{enumerate}
\def\labelenumi{\alph{enumi})}
\setcounter{enumi}{1}
\tightlist
\item
  由 a) 推出：若 f 定义在 I × E 上，且存在一个函数 h，有定义、连续、递增且 \textgreater{} 0 于 {[}0, +∞{[} 上，使得 \(\int_{0}^{+\infty} dz / h(z) = +\infty\) ，并且 \(\|\mathbf{f}(t, \mathbf{x})\| \leqslant h(\|x\|)\) ，则 \(\mathbf{x}' = \mathbf{f}(t, \mathbf{x})\) 的每个积分都在整个 I 上有定义。
\end{enumerate}

* 5) 设 \(E\) 是实数序列 \(\mathbf{x} = (x_{n})\) （满足 \(\lim_{n\to \infty}x_n = 0\) ）所成的完备赋范空间，赋以范数 \(\| \mathbf{x}\| = \sup_n|x_n|\) 。对每个整数 \(n\geqslant 0\) ，设 \(\mathbf{e}_n\) 是除下标 \(n\) 的项等于 1 外其余各项都为零的序列；空间 \(E\) 是二维子空间 \(V_{n}\) （由 \(\mathbf{e}_n\) 与 \(\mathbf{e}_{n + 1}\) 生成）与闭子空间 \(W_{n}\) （由那些 \(\mathbf{e}_k\) ——下标异于 \(n\) 与 \(n + 1\) 者——生成）的直和。设 \(\mathbf{f}_n\) 是 \(E\) 上的连续利普希茨函数，取值于 \(E\) ，在模 \(W_{n}\) 的每个等价类上为常数，等于 \(\mathbf{e}_{n + 1} - \mathbf{e}_n\) 于连接 \(\mathbf{e}_n\) 与 \(\mathbf{e}_{n + 1}\) 的直线上，并且在 \(V_{n}\) 与球 \(\| \mathbf{x}\| \leqslant \frac{1}{4}\) 的交上等于 0。另一方面，对每个整数 \(n > 0\) ，设 \(\varphi_{n}\) 是一个实函数，在区间 \(\left[\frac{1}{n + 1},\frac{1}{n}\right]\) 上有定义且连续，在该区间的端点处等于 0，并且使得 \(\int_{1 / n + 1}^{1 / n}\varphi_n(t)dt = 1\) 。设 \(I\) 是区间 {[}0, 1{]} （在 \(\mathbf{R}\) 中）；在 \(I\times E\) 上考虑函数 \(\mathbf{f}\) ，它取值于 \(E\) ，定义如下：\(\mathbf{f}(0,\mathbf{x}) = 0\) （对每个 \(\mathbf{x}\in E\) ）；对 \(\frac{1}{n + 1}\leqslant t\leqslant \frac{1}{n}\) ，令 \(\mathbf{f}(t,\mathbf{x}) = \varphi_n(t)\mathbf{f}_n(\mathbf{x})\) 。证明 \(\mathbf{f}\) 在 \(I\times E\) 上连续且局部利普希茨，但存在一个积分 \(\mathbf{u}\) ——方程 \(\mathbf{x}' = \mathbf{f}(t,\mathbf{x})\) 的积分——它在 {[}0, 1{]} 上有定义且有界，等于 \(\mathbf{e}_n\) （当 \(t = 1 / n\) ），从而当 \(t\) 趋于 0 时没有极限。

* 6) a) 设 \(\mathrm{I} = [0, +\infty[\) ；若 \(\mathbf{f}\) 在 \(\mathrm{I} \times \mathrm{E}\) 上以一个规则函数 \(k(t) > 0\) 为利普希茨函数，且积分 \(\int_0^{+\infty} k(t) dt\) 收敛，证明方程 \(\mathbf{x}' = \mathbf{f}(t, \mathbf{x})\) 的每个积分都在整个 \(\mathrm{I}\) 上有定义。

\begin{enumerate}
\def\labelenumi{\alph{enumi})}
\setcounter{enumi}{1}
\tightlist
\item
  进一步，若积分 \(\int_{0}^{+\infty}\|\mathbf{f}(t,\mathbf{x}_{0})\|dt\) （对某一点 \(x_{0}\in E\) ）收敛，证明 \(\mathbf{x}^{\prime}=\mathbf{f}(t,\mathbf{x})\) 的每个积分当 t 趋于 \(+\infty\) 时趋于一个有限极限（先证明每个积分当 t 趋于 \(+\infty\) 时有界）。
\end{enumerate}

\begin{enumerate}
\def\labelenumi{\arabic{enumi})}
\setcounter{enumi}{6}
\tightlist
\item
  考虑数量微分方程组
\end{enumerate}

\[
\frac {d x _ {i}}{d t} = \sum_ {j, k = 1} ^ {n} c _ {i j k} x _ {j} x _ {k} \quad (1 \leqslant i \leqslant n)
\]

其中 \(c_{ijk}\) 是满足 \(c_{kji} = -c_{ijk}\) 的常数。证明这个方程组的积分对 t 的一切值都有定义（注意 \(\sum_{i=1}^{n} x_{i}^{2}\) 对每个积分 \(\mathbf{x} = (x_{i})\) 都是常数）。

\begin{enumerate}
\def\labelenumi{\arabic{enumi})}
\setcounter{enumi}{7}
\tightlist
\item
  设 \(\mathbf{f}\) 是定义在 \(\mathrm{I} \times \mathrm{H}\) 上、满足 IV, p.~165 的引理 1 的条件的函数，并且对常数 \(k\) （ \(0 < k < 1\) ）与点 \(t_0 \in \mathrm{I}\) ，对所有的 \(t \in \mathrm{I}\) 与 \(\mathbf{x}_1\) 、\(\mathbf{x}_2 \in \mathrm{H}\) 有
\end{enumerate}

\[
\left\| \mathbf {f} (t, \mathbf {x} _ {1}) - \mathbf {f} (t, \mathbf {x} _ {2}) \right\| \leqslant \frac {k}{\left| t - t _ {0} \right|} \left\| \mathbf {x} _ {1} - \mathbf {x} _ {2} \right\|.
\]

证明：若 \(\mathbf{u}\) 与 \(\mathbf{v}\) 是方程 \(\mathbf{x}' = \mathbf{f}(t, \mathbf{x})\) 的两个近似解，精度分别为 \(\varepsilon_1\) 与 \(\varepsilon_2\) ，都定义在 I 上，取值于 H，并且在点 \(t_0\) 处取相同的值，则对每个 \(t \in I\) 有

\[
\left\| \mathbf {u} (t) - \mathbf {v} (t) \right\| \leqslant \frac {\varepsilon_ {1} - \varepsilon_ {2}}{1 - k} | t - t _ {0} |.
\]

由此推出，在所指出的条件下，定理 1 (IV, p.~171) 与定理 2 (IV, p.~172) 对方程 \(\mathbf{x}' = \mathbf{f}(t, \mathbf{x})\) 仍然成立。

\begin{enumerate}
\def\labelenumi{* \arabic{enumi})}
\setcounter{enumi}{8}
\tightlist
\item
  设 I 是 \(\mathbf{R}\) 中的一个区间，\(t_0\) 是 I 的一点，S 是 E 中半径 \(r\) 、中心 \(\mathbf{x}_0\) 的球，G 是由 \(\mathrm{I} \times \mathrm{S}\) 到 E 中的有界映射所成的赋范空间，其中一个映射 \(\mathbf{f}\) 的范数 \(\| \mathbf{f} \|\) 是 \(\| \mathbf{f}(t, \mathbf{x}) \|\) 在 \(\mathrm{I} \times \mathrm{S}\) 上的上确界。对每个 \(M > 0\) ，设 \(\mathrm{G}_{\mathrm{M}}\) 是 G 中的球 \(\| \mathbf{f} \| \leqslant \mathrm{M}\) 。设 L 是 G 中由 \(\mathrm{I} \times \mathrm{S}\) 到 E 中的利普希茨映射所成的子集；对每个函数 \(\mathbf{f} \in \mathrm{L} \cap \mathrm{G}_{\mathrm{M}}\) ，设 \(\mathbf{u} = \mathrm{U}(\mathbf{f})\) 是方程 \(\mathbf{x}' = \mathbf{f}(t, \mathbf{x})\) 的满足 \(\mathbf{u}(t_0) = \mathbf{x}_0\) 的积分，它取值于 S，且定义在交 \(\mathrm{J}_{\mathrm{M}}\) 上，这个交是 I 与区间
\end{enumerate}

\[
\left] t _ {0} - \frac {r}{\mathbf {M}}, t _ {0} + \frac {r}{\mathbf {M}} \right[ \text {(th. 1).}
\]

\begin{enumerate}
\def\labelenumi{\alph{enumi})}
\tightlist
\item
  设 \((\mathbf{f}_{n})\) 是属于 \(L \cap G_{M}\) 的函数序列；若 \(f_{n}\) 在 \(I \times S\) 上一致收敛于函数 f，则序列 \(\mathbf{u}_{n} = \mathrm{U}(\mathbf{f}_{n})\) 在由 \(J_{M}\) 到 E 中的有界映射所成的空间 F 中的每个聚点（关于紧收敛拓扑）都是 \(\mathbf{x}' = \mathbf{f}(t, \mathbf{x})\) 的取值 \(x_{0}\) 于点 \(t_{0}\) 的积分。反之，\(\mathbf{x}' = \mathbf{f}(t, \mathbf{x})\) 的每个定义在 \(J_{M}\) 上且满足 \(\mathbf{v}(t_{0}) = \mathbf{x}_{0}\) 的积分 v，也是某个方程 \(\mathbf{x}' = \mathbf{g}(t, \mathbf{x})\) 的积分，其中 g 是利普希茨的并且在 G 中任意接近于 f（考虑方程
\end{enumerate}

\[
\mathbf {x} ^ {\prime} = \mathbf {f} _ {n} (t, \mathbf {x}) + \mathbf {v} ^ {\prime} (t) - \mathbf {f} _ {n} (t, \mathbf {v} (t))).
\]

\begin{enumerate}
\def\labelenumi{\alph{enumi})}
\setcounter{enumi}{1}
\tightlist
\item
  进一步假设 E 是有限维的。证明：若 \(f \in G_{M}\) 满足引理 1 的假设，则对每个 \(t \in J_{M}\) ，\(H(t)\) ——即 \(\mathbf{x}' = \mathbf{f}(t, \mathbf{x})\) 的取值 \(x_{0}\) 于点 \(t_{0}\) 的那些积分在点 t 处的值所成之集——是紧连通集（要看出 \(H(t)\) 是闭的，用阿斯科利定理；要看出它是连通的，用 a)：若 \(x_{1}, x_{2}\) 是 \(H(t)\) 的两点且 \(\varepsilon > 0\) 任意，证明存在一个连通集 \(\Phi\) （由属于 \(L \cap G_{M}\) 的函数 g 组成），使得 \(\|f - g\| \leqslant \varepsilon\) 对每个函数 \(g \in \Phi\) 成立，并且诸函数 \(U(\mathbf{g})\) 在点 t 处的值所成之集是连通的且包含 \(x_{1}\) 与 \(x_{2}\) 。最后沿一个比 \(R_{+}\) 中 0 的邻域滤子更细的超滤子过渡到极限来完成证明）。
\end{enumerate}

\begin{enumerate}
\def\labelenumi{\arabic{enumi})}
\setcounter{enumi}{9}
\tightlist
\item
  设 \(f\) 是定义在长方体 \(\mathrm{P}: |t - t_0| < a\) , \(|x - x_0| < b\) （属于 \(\mathbf{R}^2\) ）上的有界连续实函数。设 \(\mathbf{M}\) 是 \(|f(t,x)|\) 在 \(\mathrm{P}\) 上的上确界，并设 \(\mathrm{I} = ]t_0 - \alpha, t_0 + \alpha[\) ，其中 \(\alpha = \inf(a, b/\mathrm{M})\) 。证明：积分集 \(\Phi\) ——即 \(x' = f(t,x)\) 的、定义在 \(\mathrm{I}\) 上且取值 \(x_0\) 于点 \(t_0\) 的积分所成之集——的上包络与下包络仍是 \(x' = f(t,x)\) 的积分，分别称为这个方程对应于点 \((t_0, x_0)\) 的最大积分与最小积分（注意，集 \(\Phi\) 对 \(\mathrm{I}\) 上的一致收敛拓扑是等度连续且闭的）。
\end{enumerate}

对每个 \(\tau \in I\) ，设 \(\xi\) 是最小积分（对应于点 \((t_0, x_0)\) 的）在点 \(\tau\) 处的值。证明：对应于点 \((\tau, \xi)\) 的最小积分与对应于点 \((t_0, x_0)\) 的最小积分在一个形如 \([\tau, \tau + h[\) 的区间上（若 \(\tau > t_0\) ）、或在一个形如 \(]\tau - h, \tau]\) 的区间上（若 \(\tau < t_0\) ）恒等。

由此推出：存在一个最大的开区间 \(]t_{1}, t_{2}[\) ，它包含 \(t_{0}\) 且含于 \(]t_{0}-a, t_{0}+a[\) ，使得对应于 \((t_{0}, x_{0})\) 的最小积分 u 可以连续地延拓到 \(]t_{1}, t_{2}[\) 上，且在 \(]t_{1}, t_{2}[\) 的每一点 t 处，值 \(u(t)\) 属于 \(]x_{0}-b, x_{0}+b[\) ，并且 u 与对应于点 \((t, u(t))\) 的最小积分在一个形如 \([t, t+h[\) 的区间上（若 \(t > t_{0}\) ）、或在一个形如 \(]t-h, t]\) 的区间上（若 \(t < t_{0}\) ）恒等；进一步证明：或者 \(t_{1} = t_{0} - a\) （相应地 \(t_{2} = t_{0} + a\) ），或者 \(\lim_{t \to t_{1}} u(t) = x_{0} \pm b\) （相应地 \(\lim_{t \to t_{2}} u(t) = x_{0} \pm b\) ）。

\begin{enumerate}
\def\labelenumi{\arabic{enumi})}
\setcounter{enumi}{10}
\tightlist
\item
  \begin{enumerate}
  \def\labelenumii{\alph{enumii})}
  \tightlist
  \item
    在长方体 P: \(|t - t_{0}| < a\) , \(|x - x_{0}| < b\) 上，设 g 与 h 是两个连续实函数，使得在 P 上 \(g(t, x) < h(t, x)\) 。设 u（相应地 v）是 \(x' = g(t, x)\) （相应地 \(x' = h(t, x)\) ）的满足 \(u(t_{0}) = x_{0}\) （相应地 \(v(t_{0}) = x_{0}\) ）的积分，定义在区间 \([t_{0}, t_{0} + c[;\) 上；证明对 \(t_{0} < t < t_{0} + c\) 有 \(u(t) < v(t)\) （考虑使这个不等式成立的 t 所成之集的上确界 \(\tau\) ）。
  \end{enumerate}
\end{enumerate}

\begin{enumerate}
\def\labelenumi{\alph{enumi})}
\setcounter{enumi}{1}
\item
  设 u 是 \(x' = g(t, x)\) 的对应于点 \((t_0, x_0)\) 的最大积分 (exerc. 10)，定义在区间 \([t_0, t_0 + c[\) 上，取值于 \(|x - x_0| < b\) 。证明：在每个紧区间 \([t_0, t_0 + d]\) （含于 \([t_0, t_0 + c[\) ）中，方程 \(x' = g(t, x) + \varepsilon\) 的对应于点 \((t_0, x_0)\) 的最小积分与最大积分只要 \(\varepsilon > 0\) 足够小就有定义，并且当 \(\varepsilon\) 沿 \textgreater{} 0 的值趋于 0 时一致收敛于 u。
\item
  设 g 与 h 是定义在 P 上的两个连续实函数，且在 P 上 \(g(t, x) \leqslant h(t, x)\) 。设 \([t_0, t_0 + c[\) 是一个区间，在其上定义了 \(x' = g(t, x)\) 的一个满足 \(u(t_0) = x_0\) 的积分 u，以及 \(x' = h(t, x)\) 的对应于点 \((t_0, x_0)\) 的最大积分 v。证明在这个区间上 \(u(t) \leqslant v(t)\) （利用 b) 化为情形 a)）。
\end{enumerate}

\begin{enumerate}
\def\labelenumi{\arabic{enumi})}
\setcounter{enumi}{11}
\tightlist
\item
  设 \(u\) 是方程 \(x' = \lambda + \frac{x^2}{1 + t^2}\) 的当 \(t = 0\) 时等于 0 的积分，并设 \(J\) 是以 0 为左端点、\(u\) 在其上连续的最大区间。
\end{enumerate}

\begin{enumerate}
\def\labelenumi{\alph{enumi})}
\item
  证明：若 \(\lambda \leqslant \frac{1}{4}\) ，则 \(J = [0, +\infty[\) （用 IV, p.~199 的 exerc. 4）。
\item
  证明：若 \(\lambda >\frac{1}{4}\) ，则 \(J = [0,a[\) ，其中
\end{enumerate}

\[
\sinh \frac {\pi}{2 \sqrt {\lambda}} <   a <   \sinh \frac {\pi}{\sqrt {4 \lambda - 1}}
\]

（令 \(x = y\sqrt{1 + t^{2}}\) 并用 exerc. 11）。

* 13) a) 设 \(\mathrm{I} = [t_0, t_0 + c]\) ，并设 \(\omega\) 是一个连续实函数，\(\geqslant 0\) ，定义在 \(\mathrm{I} \times \mathbf{R}\) 上。设 \(\mathrm{S}\) 是一个以 \(\mathbf{x}_0\) 为中心的球，位于一个完备赋范空间 \(\mathrm{E}\) 中，并设 \(\mathbf{f}\) 是一个由 \(\mathrm{I} \times \mathrm{S}\) 到 \(\mathrm{E}\) 中的连续映射，使得对任意 \(t \in \mathrm{I}\) 、\(\mathbf{x}_1 \in \mathrm{S}\) 与 \(\mathbf{x}_2 \in \mathrm{S}\) 有 \(\| \mathbf{f}(t, \mathbf{x}_1) - \mathbf{f}(t, \mathbf{x}_2) \| \leqslant \omega(t, \| \mathbf{x}_1 - \mathbf{x}_2 \|)\) 。设 \(\mathbf{u}\) 与 \(\mathbf{v}\) 是 \(\mathbf{x}' = \mathbf{f}(t, \mathbf{x})\) 的两个积分，定义在 \(\mathrm{I}\) 上，取值于 \(\mathrm{S}\) ，且满足 \(\mathbf{u}(t_0) = \mathbf{x}_1\) 、\(\mathbf{v}(t_0) = \mathbf{x}_2\) ；设 \(w\) 是 \(z' = \omega(t, z)\) 的对应于点 ( \(t_0, \| \mathbf{x}_1 - \mathbf{x}_2 \|\) ) 的最大积分 (exerc. 10)，假设它定义在 \(\mathrm{I}\) 上；证明在 \(\mathrm{I}\) 上有 \(\| \mathbf{u}(t) - \mathbf{v}(t) \| \leqslant w(t)\) 。（设 \(w(t, \varepsilon)\) 是 \(z' = \omega(t, z) + \varepsilon\) 的对应于点 ( \(t_0, \| \mathbf{x}_1 - \mathbf{x}_2 \|\) ) 的最大积分，其中 \(\varepsilon > 0\) 足够小。证明 \(\| \mathbf{u}(t) - \mathbf{v}(t) \| \leqslant w(t, \varepsilon)\) 对每个 \(\varepsilon > 0\) 成立（用反证法）。）

\begin{enumerate}
\def\labelenumi{\alph{enumi})}
\setcounter{enumi}{1}
\tightlist
\item
  在 a) 的假设陈述中，把 I 换成区间 \(I' = ]t_{0} - c, t_{0}]\) 。证明：若 w 是 \(z' = \omega(t, z)\) 的对应于点 \((t_{0}, \| \mathbf{x}_{1} - \mathbf{x}_{2}\|)\) 的在这个区间上的最小积分，则 \(\| \mathbf{u}(t) - \mathbf{v}(t) \| \geqslant w(t)\) 在 \(I'\) 上成立（同样方法）。
\end{enumerate}

* 14) a) 设 \(\omega(t,z)\) 是一个连续实函数，\(\geqslant 0\) ，对 \(0 < t \leqslant a\) 与 \(z \geqslant 0\) 有定义。假设 \(z = 0\) 是 \(z' = \omega(t,z)\) 的唯一积分，它对 \(0 < t \leqslant a\) 有定义，且满足 \(\lim_{t \to 0} z(t) = 0\) 与 \(\lim_{t \to 0} z(t)/t = 0\) 。设 \(\mathrm{I} = [\mathrm{t}_0, t_0 + a[\) ，设 \(\mathrm{S}\) 是一个以 \(\mathbf{x}_0\) 为中心的球（在 \(\mathrm{E}\) 中），\(\mathbf{f}\) 是一个由 \(\mathrm{I} \times \mathrm{S}\) 到 \(\mathrm{E}\) 中的映射，使得对任意 \(t \in \mathrm{I}\) 、\(\mathbf{x}_1 \in \mathrm{S}\) 与 \(\mathbf{x}_2 \in \mathrm{S}\) 有 \(\|\mathbf{f}(t, \mathbf{x}_1) - \mathbf{f}(t, \mathbf{x}_2)\| \leqslant \omega(|t - t_0|, \| \mathbf{x}_1 - \mathbf{x}_2 \|)\) 。证明：在一个以 \(t_0\) 为左端点、含于 \(\mathrm{I}\) 的区间上，方程 \(\mathbf{x}' = \mathbf{f}(t, \mathbf{x})\) 只能有一个解 \(\mathbf{u}\) 满足 \(\mathbf{u}(t_0) = \mathbf{x}_0\) 。（用反证法：若 \(\mathbf{v}\) 是 \(\mathbf{x}' = \mathbf{f}(t, \mathbf{x})\) 的第二个积分，从下方估计 \(\|\mathbf{u}(t) - \mathbf{v}(t)\|\) 于一个以 \(t_0\) 为左端点的区间上（借助 exerc. 13 b)），从而得出矛盾。）

应用于 \(\omega(t, z) = k(z/t)\) （其中 \(0 \leqslant k < 1\) ）的情形 (cf.~IV, p.~200, exerc. 8)。

\begin{enumerate}
\def\labelenumi{\alph{enumi})}
\setcounter{enumi}{1}
\item
  a) 的结果适用于 \(\omega(t,z)=z/t\) ；但在这种情形用一个例子证明：若 u、v 是以 \(\varepsilon\) 为精度的 \(\mathbf{x}^{\prime}=\mathbf{f}(t,\mathbf{x})\) 的两个近似积分，取值 \(x_{0}\) 于点 \(t_{0}\) ，则不可能用一个只依赖于 t（而不依赖于函数 f）的数从上方控制 \(\|\mathbf{u}(t)-\mathbf{v}(t)\|\) 。（取 f 为一个由 \(R_{+}\times R\) 到 R 中的连续映射，它对 \(t\geqslant\alpha\) 以及对 \(0<t<\alpha\) 且 \(|x|\leqslant t^{2}/(\alpha-t)\) 等于 x/t，而在其余的点 \((t,x)\) 处与 x 无关。）
\item
  设 \(\theta\) 是一个连续实函数，\(\geqslant 0\) 于区间 \([0, a]\) 上。证明：若积分 \(\int_0^a \frac{\theta(t)}{t} dt\) 收敛，则 \(a)\) 的结果适用于 \(\omega(t, z) = \frac{1 + \theta(t)}{t} z\) ；另一方面，若 \(\int_0^a \frac{\theta(t)}{t} dt\) 为无穷，给出一个连续函数 \(\mathbf{f}\) 的例子，使得 \(\| \mathbf{f}(t, \mathbf{x}_1) - \mathbf{f}(t, \mathbf{x}_2) \| \leqslant \frac{1 + \theta(|t - t_0|)}{|t - t_0|} \| \mathbf{x}_1 - \mathbf{x}_2 \|\) ，但方程 \(\mathbf{x}' = \mathbf{f}(t, \mathbf{x})\) 有无穷多个积分等于 \(\mathbf{x}_0\) 于点 \(t_0\)（方法与 \(b\) 的类似）。
\end{enumerate}

\begin{enumerate}
\def\labelenumi{\arabic{enumi})}
\setcounter{enumi}{14}
\item
  设 \(f\) 是一个实函数，对 \(|t| \leqslant a\) , \(|x| \leqslant b\) 有定义且连续，使得 \(f(t,x) < 0\) （对 \(tx > 0\) ）且 \(f(t,x) > 0\) （对 \(tx < 0\) ）；证明 \(x = 0\) 是方程 \(x' = f(t,x)\) 的在点 \(t = 0\) 处取值 0 的唯一积分（用反证法）。
\item
  设 E 是一个有限维向量空间，f 是 I×H 上的有界函数，满足 IV, p.~165 的引理 1 的条件，使得方程 \(\mathbf{x}^{\prime}=\mathbf{f}(t,\mathbf{x})\) 有唯一的解 u，它定义在 I 上，取值于 H，等于 \(x_{0}\) 于点 \(t_{0}\) 。进一步假设：对足够大的 n，存在一个近似积分 \(u_{n}\) ——\(\mathbf{x}^{\prime}=\mathbf{f}(t,\mathbf{x})\) 的、精度为 1/n 的积分——定义在 I 上，取值于 H，等于 \(x_{0}\) 于点 \(t_{0}\) 。证明序列 \((\mathbf{u}_{n})\) 在含于 I 的每个紧区间上一致收敛于 u（利用序列 \((\mathbf{u}_{n})\) 在 I 上等度连续这一事实）。
\item
  为研究方程 \(x' = \sin tx\) (IV, p.~174, 例 4)，令 u = xy，并考虑相应的方程 \(u' = \frac{u}{t} + t \sin u = \mathrm{F}(t, u)\) 的为 \textgreater0 （对 t \textgreater0 ）的解（对每个这种解有 \(u(0) = 0\) ）。以 \(\Gamma_k\) 表示由关系 \((2k - 1)\pi < u < 2k\pi\) , \(t \sin u + \frac{u}{t} = 0\) （对每个整数 \(k \geqslant 1\) ）所定义的曲线；以 \(D_k\) 表示由 \((2k - 1)\pi < u < 2k\pi\) , \(t \sin u + \frac{u}{t} < 0\) 定义的开集；以 E 表示诸 \(\overline{D}_k\) 的并在 \((t, u)\) （满足 t \textgreater0 且 u \textgreater0）所成之集中的余集。
\end{enumerate}

\begin{enumerate}
\def\labelenumi{\alph{enumi})}
\item
  证明：每条与直线 \(u = 2k\pi\) 相交的积分曲线也与直线 \(u = (2k + 1)\pi\) 相交。
\item
  若积分曲线 C 与曲线 \(\Gamma_{k}\) 交于一点 \((t_{0}, u_{0})\) ，则函数 u 对 0 \textless{} t \textless{} t \textless{} \(t_{0}\) 递增，对 \(t > t_{0}\) 递减，且当 t 趋于 \(+\infty\) 时 \(u(t)\) 趋于 \((2k - 1)\pi\) ，而曲线 C 保持含于 \(D_{k}\) 中（对 \(t > t_{0}\) ）。
\item
  证明：不存在包含在 E 中且使得 \(u(t)\) 当 t 趋于 \(+\infty\) 时趋于 \(+\infty\) 的积分曲线 C。（建立沿 C 的 x 与 u 之间的微分方程 dx/du = G(u, x)。若 C 整个位于 E 中，则 \(x^{2} > u\) （对 \(u = (2k - \frac{1}{2})\pi\) ）；另一方面，利用前面的微分方程估计 \(x^{2}(u)\) ，并得出矛盾。）
\end{enumerate}

\(d\)) 证明：对每个整数 \(k\) ，存在一条且仅一条包含在 E 中且使得 \(u(t)\) 趋于 \(2k\pi\) （当 \(t\) 趋于 \(+\infty\) 时）的积分曲线 C。（令 \(v = 2k\pi - u\) ，得到方程 \(v' = \frac{v - 2k\pi}{t} + t \sin v\) ；把该方程与 \(v' = \frac{v - 2k\pi}{t} + tv\) 相比较，其中 \(t\) 趋于 \(+\infty\) 而 \(v\) 趋于 0。）

\begin{enumerate}
\def\labelenumi{\arabic{enumi})}
\setcounter{enumi}{17}
\tightlist
\item
  设 E 是实数序列 \(\mathbf{x}=(x_{n})_{n\in\mathbb{N}}\) （满足 \(\lim_{n\to\infty}x_{n}=0\) ）所成的空间，赋以范数 \(\|x\|=\sup_{n}|x_{n}|\) ，它实际上是一个完备赋范空间。对每个 \(\mathbf{x}=(x_{n})\in\mathrm{E}\) ，设 \(\mathbf{y}=\mathbf{f}(\mathbf{x})\) 是 E 的元素 \((y_{n})\) ，由 \(y_{n}=|x_{n}|^{\frac{1}{2}}+\frac{1}{n+1}\) 定义；函数 f 在 E 上连续。证明：微分方程 \(\mathbf{x}^{\prime}=\mathbf{f}(\mathbf{x})\) 没有定义在 R 中 0 的一个邻域上、取值于 E 且在点 t=0 处等于 0 的解 (cf.~IV, p.~202, exerc. 11)。
\end{enumerate}

\setcounter{exercisegroup}{1}
\refstepcounter{exercisegroup}
